\documentclass[11pt]{article}
\usepackage[T1]{fontenc}
\usepackage[utf8]{inputenc}
\usepackage{lmodern}
\usepackage[margin=1in]{geometry}
\usepackage{amsmath,amssymb,amsthm}
\usepackage{microtype}
\usepackage[colorlinks=true,linkcolor=blue,citecolor=blue,urlcolor=blue]{hyperref}
\newtheorem{theorem}{Theorem}
\newtheorem{lemma}{Lemma}
\newcommand{\one}{\mathbf 1}
\newcommand{\diag}{\operatorname{diag}}
\newcommand{\osc}{\operatorname{osc}}
\newcommand{\relint}{\operatorname{relint}}
\newcommand{\Tan}{\operatorname{Tan}}
\title{Interior smoothness of carrying simplices\protect\\
for competitive population maps}
\author{Kaiwen Jin\protect\\
\small\href{mailto:yui.ui.siki@gmail.com}{yui.ui.siki@gmail.com}}
\date{October 12, 2026}
\begin{document}
\maketitle

\begin{abstract}
A carrying simplex reduces the nonzero long-term dynamics of a competitive
population map to an invariant hypersurface. We prove that its positive
interior is a $C^1$ embedded manifold under the standard $C^1$ per-capita
competition and spectral-radius assumptions. Equivalently, its radial
function is $C^1$ in the interior of the probability simplex.
The proof uses logarithmic coordinates. Although inverse derivatives lose
strict positivity at extinction faces, their weighted columns give a uniform
contraction on a fixed cone of positive vectors. The resulting continuous
covector annihilates every contingent direction of the logarithmic simplex
and forces its Lipschitz graph to be $C^1$.
No boundary smoothness, convexity or permanence assumption is imposed.
\end{abstract}

\section{The regularity question}

Consider a discrete population model on $C=[0,\infty)^d$ of the form
$F_i(x)=x_i f_i(x)$. The factor $f_i$ describes per-capita growth. Strict
decrease of every $f_i$ with respect to every population density expresses
both interspecific and intraspecific competition. A carrying simplex is a
compact invariant set $\Sigma\subset C\setminus\{0\}$ which meets each
nonzero ray once, is unordered for the coordinate order, and carries the
asymptotic nonzero dynamics. Unordered means that distinct points cannot
satisfy $x_i\leq y_i$ for all $i$. Write
\begin{equation}\label{eq:radial}
 \Sigma=\{R(u)u:u\in\Delta\},\qquad
 \Delta=\{u\in C:\textstyle\sum_i u_i=1\},
\end{equation}
where $R$ is continuous and positive. This representation gives a geometric
reduction by one dimension. Smoothness is needed to use differential
geometric methods directly on the reduced dynamics.

Baigent's existence theorem \cite[Theorem~1]{Baigent} applies under the
hypotheses stated below. The smoothness of the entire positive interior is
asked in \cite[Section~6]{Baigent}. Mierczy\'nski, Niu and Ruiz-Herrera
\cite[Section~1, Corollaries~4.8--4.9 and Theorem~4.12]{MNR} study local
linearization and invariant manifolds near interior fixed points. Their local
conclusions do not assert smoothness at every interior point. Under additional
assumptions, convexity yields a $C^1$ carrying simplex neatly embedded in the
orthant \cite[Main Theorem]{Mierczynski}. That includes a boundary conclusion
which is distinct from the question considered here. Logarithmic carrying
simplices have also been studied for planar Leslie--Gower and Ricker maps,
with conditions ensuring convexity \cite{NYB}.

We answer the interior question for the full following class of maps.

\begin{theorem}\label{thm:main}
Let $d\geq3$ and $C=[0,\infty)^d$. Let $f$ be $C^1$ on an open
neighborhood of $C$, and put $F_i(x)=x_i f_i(x)$. Suppose that
\[
 f_i(x)>0,\qquad \partial_j f_i(x)<0
 \quad(x\in C,\ 1\leq i,j\leq d).
\]
For each $i$ suppose that $q_i e_i$ is a fixed point, with $q_i>0$.
Put $q=(q_1,\ldots,q_d)$ and assume
\begin{equation}\label{eq:spectral}
 \rho(M(x))<1\quad(x\in[0,q]),\qquad
 M_{ij}(x)=-\frac{x_i}{f_i(x)}\partial_j f_i(x),
 \qquad [0,q]=\prod_i[0,q_i],
\end{equation}
where $\rho$ denotes spectral radius.
Then the guaranteed carrying simplex has
$\Sigma^\circ=\Sigma\cap(0,\infty)^d$ a $C^1$ embedded
$(d-1)$-dimensional submanifold of $(0,\infty)^d$.
In \eqref{eq:radial}, $R\in C^1(\relint\Delta)$.
\end{theorem}

The theorem imposes no assumption on regularity or transverse embedding at
extinction faces. In particular, it does not assert neat embedding of the
whole simplex. The constants used in its proof may depend on the fixed map
and dimension.

The main difficulty is that backward orbits of interior points may approach
extinction faces. In physical coordinates some inverse derivative entries
then vanish, so a lower bound on every population is unavailable. In
logarithmic coordinates the inverse derivatives instead have the form
$I+B(x)\diag(x)$, with all entries of $B$ bounded above and below by positive
constants. Compactness separates the simplex from the origin, giving a lower
bound on the \emph{sum} of the populations. These facts yield uniform
contraction of ratios on a fixed positive cone. A second normalization gives
a continuous covector whose kernel contains all contingent directions.
Unorderedness then turns this constraint into differentiability of the
logarithmic graph. The proof below gives both normalizations explicitly.

\section{Proof of Theorem~\ref{thm:main}}

All vector inequalities in this section are coordinatewise. Logarithms and
exponentials are applied coordinatewise. For a vector $r$ write
$\osc(r)=\max_i r_i-\min_i r_i$.

\subsection{The logarithmic graph and its inverse derivatives}

Each axial fixed point $q_i e_i$ belongs to $\Sigma$. Baigent's
Theorem~1 \cite{Baigent} gives distance attraction:
$\operatorname{dist}(F^n x,\Sigma)\to0$ for every $x\in C\setminus\{0\}$.
For $x=q_i e_i$ this distance is constant, so it is zero; closedness of
$\Sigma$ gives membership. If $x\in\Sigma$ had $x_i>q_i$, it would be
comparable with $q_i e_i$. Therefore
\begin{equation}\label{eq:mass}
 \Sigma\subset[0,q],\qquad
 s:=\min_{x\in\Sigma}\sum_i x_i>0,\qquad L:=\sum_i q_i.
\end{equation}
The minimum is positive because $\Sigma$ is compact and excludes zero.

Under \eqref{eq:spectral} and the sign assumptions, Baigent's
Lemma~16(1) \cite{Baigent} states that $F|_{[0,q]}$ is a homeomorphism
onto its image. The source hypotheses AS3--AS4 are exactly the positivity,
derivative signs and spectral condition here. The axial point on each axis
is unique, since $t\mapsto f_i(t e_i)$ is strictly decreasing.
Since $F(\Sigma)=\Sigma$, we may set
\[
 T=(F|_\Sigma)^{-1}.
\]
This is a continuous map on the compact physical simplex. Positivity of
$f_i$ implies that $F$ preserves support, so $T$ preserves $\Sigma^\circ$.

Put $\one=(1,\ldots,1)$ and
\[
 E=\{w\in\mathbb R^d:\textstyle\sum_i w_i=0\},\qquad
 \pi z=z-\frac1d\Big(\sum_i z_i\Big)\one,
 \qquad S=\log\Sigma^\circ.
\]
The line $w+\mathbb R\one$ exponentiates to a positive ray. By the radial
representation it meets $S$ exactly once, so there is a function $h:E\to
\mathbb R$ with
\begin{equation}\label{eq:graph}
 S=\{w+h(w)\one:w\in E\}.
\end{equation}
The exponential preserves order, and hence $S$ is unordered. For $w,v\in E$
this gives
\[
 \min_i(w_i-v_i+h(w)-h(v))\leq0\leq
 \max_i(w_i-v_i+h(w)-h(v)).
\]
Consequently
\begin{equation}\label{eq:lipschitz}
 |h(w)-h(v)|\leq\|w-v\|_\infty\leq\|w-v\|_2.
\end{equation}

We next record bounds for the logarithmic inverse derivative. Define on
the closed box
\begin{equation}\label{eq:matrices}
 K_{ij}(x)=-\frac{\partial_j f_i(x)}{f_i(x)},\qquad
 N(x)=K(x)\diag(x),\qquad A(x)=(I-N(x))^{-1}.
\end{equation}
All entries of $K$ are positive. The identity
$\det(I-tUV)=\det(I-tVU)$, obtained by the two block eliminations of
$\left(\begin{smallmatrix}I&tU\\V&I\end{smallmatrix}\right)$,
shows that $N$ and $M=\diag(x)K$ have identical characteristic polynomials,
including at the boundary. Thus $\rho(N)<1$ and $A$ is continuous on the
box. The nonnegative matrix expansion in the proof of
\cite[Lemma~15]{Baigent} gives $(I-M)^{-1}\geq I$.
At a point with all coordinates positive,
$N=\diag(x)^{-1}M\diag(x)$, so $A\geq I$ there. Approximating a boundary
point by $(1-\varepsilon)x+\varepsilon q/2$ and using continuity gives
$A\geq I$ everywhere on the box.
Set $B=AK$. Then
\begin{equation}\label{eq:columns}
 A(x)=I+B(x)\diag(x),\qquad
 0<a\leq B_{ij}(x)\leq b<\infty
 \quad(x\in[0,q],\ 1\leq i,j\leq d).
\end{equation}
Such $a,b$ exist because $B\geq K>0$ is continuous on a compact set.
Choose
\begin{equation}\label{eq:constants}
 \kappa=\frac{a}{2b},\qquad
 \eta=\frac{a\kappa s}{1+bL},\qquad \tau=1-\eta.
\end{equation}
Here $0<\kappa\leq1/2$ and $0<\eta<1$, so $0<\tau<1$.

For the logarithmic map
$H(z)=\log F(e^z)=z+\log f(e^z)$ we have
$DH(\log x)=I-N(x)$. The inverse function theorem provides a local $C^1$
inverse at every positive point of $\Sigma$. Continuity of $T$ identifies
these branches, on nearby points of $S$, with
$G(\log x)=\log Tx$. Thus, for every fixed positive integer $n$,
\begin{equation}\label{eq:inverse-product}
 DG(\log x)=A(Tx),\qquad
 DG^n(\log x)=P_n(x):=A(T^nx)\cdots A(Tx).
\end{equation}
In this formula $G^n$ is a local composition along a finite orbit segment.

\subsection{Uniform contraction of positive ratios}

Consider the cone
\[
 \mathcal C_\kappa=\{v\in(0,\infty)^d:
               \min_i v_i\geq\kappa\max_i v_i\}.
\]
It keeps coordinate ratios bounded. If $v\in\mathcal C_\kappa$ and
$c=\sum_j x_jv_j$, \eqref{eq:columns} gives
\[
 \min_i(A(x)v)_i\geq\min_i v_i+ac,
 \qquad
 \max_i(A(x)v)_i\leq\max_i v_i+bc.
\]
Since $a-\kappa b\geq0$, $A(x)$ preserves this cone for every point of
the box. On positive rays use the Hilbert distance
\[
 d_H(v,w)=\osc\big((\log(v_i/w_i))_{i=1}^d\big).
\]
It is unchanged by positive scalar multiples. Its diameter on
$\mathcal C_\kappa$ is at most $D:=2\log(1/\kappa)$.

\begin{lemma}\label{lem:contraction}
For every $x\in\Sigma$ and $v,w\in\mathcal C_\kappa$,
\begin{equation}\label{eq:contraction}
 d_H(A(x)v,A(x)w)\leq\tau d_H(v,w).
\end{equation}
\end{lemma}
\begin{proof}
The Jacobian of $r\mapsto\log(A(x)e^r)$ at $e^r=v$ has entries
\[
 J_{ij}=\frac{A_{ij}(x)v_j}{(A(x)v)_i}.
\]
It is row stochastic, meaning that its entries are nonnegative and each
row sums to one. From \eqref{eq:columns} and $v\in\mathcal C_\kappa$,
\begin{equation}\label{eq:common-row}
 J_{ij}\geq\nu_j(x):=\frac{a\kappa x_j}{1+bL},\qquad
 \sum_j\nu_j(x)\geq\eta.
\end{equation}
Indeed, $A_{ij}(x)\geq ax_j$, $v_j\geq\kappa\max_kv_k$, and
$(A(x)v)_i\leq(1+bL)\max_kv_k$.
If $m=\sum_j\nu_j(x)$, all rows of $J-\one\nu$ have nonnegative
entries and mass $1-m$. Subtracting the common scalar $\nu r$ gives
\begin{equation}\label{eq:osc}
 \osc(Jr)\leq(1-m)\osc(r)\leq\tau\osc(r).
\end{equation}
Here $m\leq a\kappa L/(1+bL)<1$.

For $v,w$ in the cone, the segment
$r(t)=(1-t)\log w+t\log v$ has $e^{r(t)}\in\mathcal C_\kappa$:
every coordinate difference is at least $\log\kappa$.
The integral mean-value formula therefore gives
\[
 \log(A(x)v)-\log(A(x)w)
 =\int_0^1 J(r(t))(\log v-\log w)\,dt.
\]
For each pair of components their difference is bounded by
$\tau\osc(\log v-\log w)$, using \eqref{eq:osc}. Taking the largest
pairwise difference after integration proves \eqref{eq:contraction}.
\end{proof}

Normalize positive vectors by their coordinate sum and write
$\widehat v=v/(\sum_i v_i)$. The normalized cone section
$V=\{v\in\mathcal C_\kappa:\sum_i v_i=1\}$ is compact and nonempty;
every coordinate is at least $\kappa/d$. For $x\in\Sigma$ let
\[
 L_n(x)=A(x)A(Fx)\cdots A(F^{n-1}x),\qquad
 V_n(x)=\{\widehat{L_n(x)v}:v\in V\}.
\]
These are nested nonempty compact subsets of $V$, since every factor preserves
the cone. Lemma~\ref{lem:contraction} bounds their Hilbert diameters by
$D\tau^n$. Their intersection is consequently a single point $v_*(x)$.
To see this in an ordinary norm, if $v,w$ are normalized and
$d_H(v,w)\leq\delta$, their coordinate ratios straddle one and lie in
$[e^{-\delta},e^\delta]$. Hence
\begin{equation}\label{eq:metric-conversion}
 \|v-w\|_1\leq e^\delta-1.
\end{equation}
The continuous functions $\widehat{L_n(x)\one}$ converge uniformly to
$v_*(x)$ with error at most $e^{D\tau^n}-1$ in this norm. Thus $v_*$
is continuous on the compact physical $\Sigma$, and
\begin{equation}\label{eq:covariance}
 \sum_i v_{*,i}(x)=1,\qquad
 v_*(x)\in\mathcal C_\kappa,\qquad
 A(Tx)v_*(x)=\gamma(x)v_*(Tx),\quad \gamma(x)>0.
\end{equation}
The last identity follows from $L_{n+1}(Tx)=A(Tx)L_n(x)$ by normalization
and passage to the limit. The fields at boundary points are auxiliary matrix
constructions; no logarithm or tangent plane is assigned there.

\subsection{A continuous covector from row products}

We normalize the inverse derivatives using the positive direction just
constructed. For $x\in\Sigma$ set
\begin{equation}\label{eq:Q}
 Q_{ij}(x)=
 \frac{A_{ij}(Tx)v_{*,j}(x)}{(A(Tx)v_*(x))_i},\qquad
 Q^{(n)}(x)=Q(T^{n-1}x)\cdots Q(Tx)Q(x).
\end{equation}
These matrices are row stochastic. Since $v_*(x)\in\mathcal C_\kappa$,
\eqref{eq:common-row} applies at $Tx$ with input $v_*(x)$: every row of
$Q(x)$ dominates $\nu(Tx)$, whose entries sum to at least $\eta$.
Thus every $Q(x)$ contracts oscillation by $\tau=1-\eta$. Therefore
\begin{equation}\label{eq:row-diameter}
 \osc(Q^{(n)}(x)r)\leq\tau^n\osc(r).
\end{equation}
The convex hulls of the rows of $Q^{(n)}(x)$ are nested, since the next
product is obtained by left multiplication by a row-stochastic matrix.
For two rows $a_1,a_2$ and $\|r\|_\infty\leq1$,
\eqref{eq:row-diameter} gives $|(a_1-a_2)r|\leq2\tau^n$.
Using the elementary duality between the $\ell^1$ and $\ell^\infty$ norms,
the entire row hull has $\ell^1$ diameter at most $2\tau^n$.
It has a unique limiting probability row $p(x)$, and all rows converge
uniformly to it. Finite products are continuous, so $p$ is continuous.

Define a row covector by
\begin{equation}\label{eq:ell}
 \ell_j(x)=\frac{p_j(x)}{v_{*,j}(x)}.
\end{equation}
It is continuous and nonnegative, and
$\ell(x)v_*(x)=1$, so $\ell(x)\one>0$. Individual entries of $p$ need
not be positive at the boundary.

For clarity, the product normalization can be written exactly. Put
$v_k=v_*(T^kx)$ and $D_k=\diag(v_k)$, and write
$A(T^kx)v_{k-1}=\gamma_k v_k$. By \eqref{eq:Q},
\[
 Q(T^{k-1}x)=\gamma_k^{-1}D_k^{-1}A(T^kx)D_{k-1}.
\]
The successive diagonal factors cancel, giving
\[
 Q^{(n)}(x)=(\gamma_1\cdots\gamma_n)^{-1}
             D_n^{-1}P_n(x)D_0,
 \qquad P_n(x)v_0=(\gamma_1\cdots\gamma_n)v_n.
\]
Thus, for every real vector $v$ and every component $i$,
\begin{equation}\label{eq:sign-limit}
 \frac{(P_n(x)v)_i}{(P_n(x)v_*(x))_i}
 =\left[Q^{(n)}(x)
       \left(\frac{v_j}{v_{*,j}(x)}\right)_{j=1}^d\right]_i
 \longrightarrow\ell(x)v.
\end{equation}
All denominators are positive. If $\ell(x)v\ne0$, then every coordinate
of $P_n(x)v$ has that same strict sign for all sufficiently large $n$.

\subsection{Contingent directions and differentiability}

For $z\in S$ let $\Tan_S(z)$ consist of limits
\[
 v=\lim_{k\to\infty}t_k(z_k-z),\qquad
 z_k\in S,\quad z_k\to z,\quad t_k>0.
\]
Unorderedness excludes strictly positive and strictly negative vectors
from this contingent cone. A sequence tending to such a vector would
eventually give points strictly comparable with $z$.

Fix $x\in\Sigma^\circ$ and $v\in\Tan_S(\log x)$. For each fixed $n$,
the local $C^1$ map $G^n$ from \eqref{eq:inverse-product} sends nearby
points of $S$ into $S$ and hence
\[
 P_n(x)v\in\Tan_S(\log T^nx).
\]
Indeed, its differentiability remainder is $o(\|z_k-\log x\|)$;
multiplication by $t_k$ still makes it tend to zero, because
$t_k\|z_k-\log x\|$ is bounded whenever the scaled increments have a
finite limit. If $\ell(x)v\ne0$, \eqref{eq:sign-limit} would make this
transported contingent vector strictly signed for some finite $n$, a
contradiction. Consequently
\begin{equation}\label{eq:tangent}
 \Tan_S(\log x)\subset\ker\ell(x).
\end{equation}
Only finite local compositions were differentiated in this argument.

Recall that $E=\{w:\sum_i w_i=0\}$ and $h:E\to\mathbb R$ is
$1$-Lipschitz by \eqref{eq:lipschitz}.
At $z=w+h(w)\one=\log x$ define the linear functional on $E$
\begin{equation}\label{eq:Dh}
 \Lambda_w(r)=-\frac{\ell(x)r}{\ell(x)\one}.
\end{equation}
Its denominator is positive. To prove that $Dh(w)=\Lambda_w$, take
any sequence $r_k\in E\setminus\{0\}$ tending to zero. By
\eqref{eq:lipschitz}, a subsequence has limits
\[
 \frac{r_k}{\|r_k\|_2}\longrightarrow r_\infty,\qquad
 \frac{h(w+r_k)-h(w)}{\|r_k\|_2}\longrightarrow c.
\]
Then $r_\infty+c\one\in\Tan_S(z)$, so \eqref{eq:tangent} forces
$c=\Lambda_w(r_\infty)$. Along this subsequence,
\[
 \frac{h(w+r_k)-h(w)-\Lambda_w(r_k)}{\|r_k\|_2}\longrightarrow0.
\]
If the full remainder limit were not zero, a sequence staying away from
zero would have precisely such a subsequence, giving a contradiction.
This proves Fr\'echet differentiability. The formula \eqref{eq:Dh},
continuity of $\ell$ and of $x=\exp(w+h(w)\one)$ show that $Dh$ is
continuous. Thus $h\in C^1(E)$.

The graph map $w\mapsto w+h(w)\one$ is injective, has continuous inverse
$\pi$ on its image, and has injective derivative because projecting that
derivative gives the identity on $E$. It is therefore a $C^1$ embedding.
Exponentiation is a $C^1$ diffeomorphism onto the positive orthant, so
$\Sigma^\circ$ is a $C^1$ embedded hypersurface.

Finally, for $u\in\relint\Delta$, \eqref{eq:graph} and the unique
intersection with its ray give
\begin{equation}\label{eq:R}
 R(u)=\exp\left(h(\pi\log u)-\frac1d\sum_i\log u_i\right).
\end{equation}
This is $C^1$ on $\relint\Delta$, where all $u_i$ are positive. The map
and dimension were arbitrary under the hypotheses of
Theorem~\ref{thm:main}, which completes its proof.\hfill$\square$

\paragraph{AI-use statement.}
OpenAI Codex assisted with literature retrieval, development of the argument,
proof checking, and preparation of the manuscript. AI review does not replace
independent mathematical review.

\begin{thebibliography}{9}
\bibitem{Baigent}
S. Baigent, \emph{Carrying Simplices for Competitive Maps}, in
\emph{Difference Equations, Discrete Dynamical Systems and Applications:
ICDEA 2017}, Springer, 2019, pp.\ 3--29.
\href{https://doi.org/10.1007/978-3-030-20016-9_1}{doi:10.1007/978-3-030-20016-9\_1}.
\href{https://homepages.ucl.ac.uk/~ucess29/resources/ICDEA_RPS.pdf}{Author manuscript}
(Theorem~1, Lemmas~15--16, and Section~6; manuscript pp.\ 9, 18 and 24).

\bibitem{MNR}
J. Mierczy\'nski, L. Niu and A. Ruiz-Herrera,
\emph{Linearization and invariant manifolds on the carrying simplex for
competitive maps}, Journal of Differential Equations \textbf{267} (2019),
7385--7410.
\href{https://doi.org/10.1016/j.jde.2019.08.001}{doi:10.1016/j.jde.2019.08.001}.
\href{https://arxiv.org/abs/1902.08914v2}{arXiv:1902.08914}.

\bibitem{Mierczynski}
J. Mierczy\'nski, \emph{The $C^1$ property of convex carrying simplices
for competitive maps}, Ergodic Theory and Dynamical Systems
\textbf{40} (2020), 1335--1350.
\href{https://doi.org/10.1017/etds.2018.85}{doi:10.1017/etds.2018.85}.
\href{https://arxiv.org/abs/1801.01032v2}{arXiv:1801.01032}.

\bibitem{NYB}
H. Naderi Yeganeh and S. Baigent,
\emph{Convexity of non-compact carrying simplices in logarithmic coordinates},
Journal of Difference Equations and Applications \textbf{30} (2024),
1671--1691.
\href{https://doi.org/10.1080/10236198.2024.2369215}{doi:10.1080/10236198.2024.2369215}.
\end{thebibliography}
\end{document}
